% Appendix: The Small Universes: D_4 and D_9
% Complete Canonical Catalogues of Sister Universes

To appreciate the vastness of the $z_4^4$ universe on $\mathbb{Z}_{16}$, it is illuminating to examine its smaller sister universes: the $D_4$ universe of 2-block partitions on $\mathbb{Z}_4$ ($z_2^2$), and the $D_9$ universe of 3-block partitions on $\mathbb{Z}_9$ ($z_3^3$). Where $\mathbb{Z}_{16}$ generates tens of thousands of orbits across multiple pages, these sibling universes are completely and exhaustively surveyable in a single glance.

\section{The \texorpdfstring{$D_4$}{D4} Universe}
\label{sec:d4-universe}

The metric arena $\mathbb{Z}_4 = \{0, 1, 2, 3\}$ represents the square. A mosaic of shape $z_2^2$ is a set partition of the four vertices into two blocks of two vertices each.

\subsection{Dihedral Orbital Enumeration}

The total number of raw equal-block partitions of shape $z_2^2$ is:
\[
\frac{1}{2!}\binom{4}{2} = \frac{6}{2} = 3
\]
These three partitions are:
\begin{align*}
P_0 &= \{\{0, 2\}, \{1, 3\}\} \\
P_1 &= \{\{0, 1\}, \{2, 3\}\} \\
P_2 &= \{\{0, 3\}, \{1, 2\}\}
\end{align*}
The dihedral group $D_4$ of order $8$ (the symmetries of the square) acts transitively on the vertices. Under this action:
\begin{itemize}
    \item $P_0$ connects opposite vertices by orthogonal interior diagonals (stride profile $(2, 2)$). It is invariant under all four rotations and all four reflections of $D_4$. Its stabiliser has full order $|\operatorname{Stab}_{D_4}(P_0)| = 8$, yielding a singleton dihedral orbit of size $|O| = 1$.
    \item $P_1$ and $P_2$ connect adjacent boundary edges (stride profile $(1, 3)$). The rotation $r \colon x \mapsto x + 1 \pmod 4$ interchanges them ($r(P_1) = P_2$, $r(P_2) = P_1$). They merge into a single dihedral orbit $\{P_1, P_2\}$ of size $|O| = 2$, with stabiliser order $4$.
\end{itemize}
Hence, there are exactly two dihedral orbits under $D_4$.

\subsection{Quotient Collapse}

The unit group of the ring $\mathbb{Z}_4$ consists of the odd residues:
\[
(\mathbb{Z}_4)^\times = \{1, 3\} = \{1, -1\}
\]
Because reflection across vertex $0$ acts on $\mathbb{Z}_4$ as $x \mapsto -x \equiv 3x \pmod 4$, the inversion multiplier $\mu_3$ is already an element of the rigid dihedral reflection group: $\mu_3 \in D_4$.

Consequently, the outer affine automorphism group collapses completely:
\[
G_{\text{macro}} = \operatorname{Aff}(\mathbb{Z}_4) / D_4 \cong (\mathbb{Z}_4)^\times / \{\pm 1\} \cong C_1
\]
The macro-quotient is the trivial group of order $1$. Every dihedral orbit is therefore an affine macro-orbit of length $1$ (a monomer). There are no dimers or tetramers.

\begin{figure}[htbp]
\centering
\input{tkz/fig_d4_universe.tkz}
\caption{\label{fig:d4-universe}The complete $D_4$ universe: the two affine monomers of shape $z_2^2$ on $\mathbb{Z}_4$. $M_1$ is the regular orthogonal cross ($|O|=1$); $M_2$ represents the parallel edges ($|O|=2$)}
\end{figure}

The complete $D_4$ catalogue consists of precisely two monomers, illustrated in Figure~\ref{fig:d4-universe}:
\begin{enumerate}
    \item \textbf{Monomer $M_1$} ($|O|=1$): The regular coset partition of the ideal $2\mathbb{Z}_4$.
    \item \textbf{Monomer $M_2$} ($|O|=2$): The parallel chord pair.
\end{enumerate}
Machine-readable relational manifests linking both monomers to \texttt{polyatlas} compositions are tabulated in \path{data/z4_mosaics.tsv} and \path{data/z4_mosaic_blocks.tsv} (see Section~\ref{sec:hexdec-relational-datasets}).

\section{The \texorpdfstring{$D_9$}{D9} Universe}
\label{sec:d9-universe}

The metric arena $\mathbb{Z}_9 = \{0, 1, \dots, 8\}$ represents the regular nonagon. A mosaic of shape $z_3^3$ partitions the nine vertices into three blocks of three vertices each (three triangles).

\subsection{The 25 Dihedral Orbits}

The raw combinatorial count of equal-block partitions of shape $z_3^3$ is:
\[
\frac{1}{3!}\binom{9}{3, 3, 3} = \frac{1}{6} \cdot \frac{9!}{(3!)^3} = \frac{1680}{6} = 280
\]
The dihedral group $D_9$ has order $|D_9| = 18$ (nine rotations and nine reflections). Under Burnside's Lemma, the number of non-equivalent dihedral orbits is:
\[
\frac{1}{18} \sum_{g \in D_9} |\operatorname{Fix}(g)| = 25
\]
These $25$ dihedral orbits stratify into four stabiliser size classes:
\begin{itemize}
    \item \textbf{Size 1 (1 orbit)}: The unique coset partition of the ideal $3\mathbb{Z}_9$:
    \[
    P_0 = \{\{0, 3, 6\}, \{1, 4, 7\}, \{2, 5, 8\}\}
    \]
    Each block is a regular equilateral triangle with stride profile $(3, 3, 3)$ (shaded blue). Invariant under all $18$ symmetries of $D_9$.
    \item \textbf{Size 3 (3 orbits)}: Three orbits with order-$6$ stabiliser $\langle r^3, s \rangle$.
    \item \textbf{Size 9 (12 orbits)}: Twelve orbits with order-$2$ bilateral reflectional stabilisers.
    \item \textbf{Size 18 (9 orbits)}: Nine fully asymmetric chiral orbits ($|\operatorname{Stab}| = 1$).
\end{itemize}
The partition sum accounts for all $280$ configurations:
\[
1 \cdot 1 + 3 \cdot 3 + 12 \cdot 9 + 9 \cdot 18 = 1 + 9 + 108 + 162 = 280
\]

\subsection{The Outer Automorphism Group \texorpdfstring{$C_3$}{C3}}

The unit group $(\mathbb{Z}_9)^\times$ consists of the six integers coprime to $9$:
\[
(\mathbb{Z}_9)^\times = \{1, 2, 4, 5, 7, 8\} \cong C_6
\]
Because reflection across vertex $0$ acts as $x \mapsto -x \equiv 8x \pmod 9$, the subgroup $\{\pm 1\} = \{1, 8\}$ lies entirely within $D_9$. Taking the quotient yields the outer macro-automorphism group:
\[
G_{\text{macro}} = \operatorname{Aff}(\mathbb{Z}_9) / D_9 \cong (\mathbb{Z}_9)^\times / \{\pm 1\} \cong C_6 / C_2 \cong C_3
\]
This outer group $G_{\text{macro}} \cong C_3$ is generated by the dilation multiplier:
\[
\mu_2 \colon x \mapsto 2x \pmod 9
\]
with cycle structure $\mu_2 \colon 1 \mapsto 2 \mapsto 4 \mapsto 8 \equiv -1 \equiv 1 \pmod{G_{\text{macro}}}$.

Because $|G_{\text{macro}}| = 3$ is a prime number, the only possible divisors are $1$ and $3$. Thus, every macro-orbit in the $D_9$ universe is either:
\begin{itemize}
    \item A \textbf{Monomer} (fixed point of length $1$).
    \item A \textbf{Trimer} (a 3-cycle under $\mu_2 \colon O \mapsto \mu_2(O) \mapsto \mu_4(O) \mapsto O$).
\end{itemize}
Applying Burnside's Lemma to the action of $C_3$ on the $25$ dihedral orbits reveals that $\mu_2$ fixes only the regular coset partition $P_0$. The remaining $25 - 1 = 24$ dihedral orbits partition into exactly $24 / 3 = 8$ Trimers!

\subsection{The Complete Inventory of Trimers}

To maintain complete methodological uniformity across universes, the nonagonal trimers are catalogued under the exact same 5-tier natural sorting cascade established for the hexadecagonal tetramers in Section~\ref{sec:tetramers-and-scale-models}:
\begin{enumerate}
    \item \textbf{Tier 1: Dihedral Orbit Size $|O|$}: Ascending order $|O| \in \{3, 9, 18\}$.
    \item \textbf{Tier 2: Symmetric Block Count $s$}: Descending order $s \in \{3, 1, 0\}$, producing the characteristic visual transition from pure cobalt blue ($s=3$) through mixed chirality ($s=1$) to full chiral pink ($s=0$).
    \item \textbf{Tier 3: Circular Chord Dispersion Energy $\mathcal{E}$}: Descending order of total chord dispersion energy.
    \item \textbf{Tier 4: Canonical Stride Profile Multiset}: Grouping by unoriented cyclic stride profiles.
    \item \textbf{Tier 5: Lexicographic Tie-Breaker}: Sorting by canonical block elements.
\end{enumerate}
Furthermore, the orientation of each 3-strip is anchored by evaluating the 5-tier key across all three cycle phases, placing the phase of highest symmetry and dispersion energy at the leading $\mu_1$ position.

The complete census of eight trimers is catalogued in Table~\ref{tab:d9-trimers}. Together with the sole monomer $M_1$, these nine cartouches constitute the entire $D_9$ universe. The visual atlas is presented in Figure~\ref{fig:d9-universe}, where each trimer is displayed as a horizontal 3-strip tracking the Galois dilation sequence $\mu_1 \to \mu_2 \to \mu_4$. Full relational block manifests connecting all $25$ dihedral comosaics to \texttt{polyatlas} are tabulated in \path{data/z9_mosaics.tsv} and \path{data/z9_mosaic_blocks.tsv}.

\begin{table}[htbp]
\centering
\footnotesize
\setlength{\tabcolsep}{4.5pt}
\begin{tabular}{cccccc}
\toprule
\textbf{Trimer} & \textbf{$|O|$} & \textbf{$s$} & \textbf{$\mathcal{E}$} & \textbf{Canonical Strides} & \textbf{Galois Dilation Cycle ($\mu_1 \to \mu_2 \to \mu_4$)} \\
\midrule
$T_1$ & 3  & 3 & 99 & $(1, 4, 4)^3$ & \texttt{/015/267/348/} \texttt{[sss]} $\to$ \texttt{/012/345/678/} $\to$ \texttt{/024/168/357/} \\
$T_2$ & 9  & 3 & 90 & $(1, 4, 4)^2, (2, 2, 5)$ & \texttt{/015/237/468/} \texttt{[sss]} $\to$ \texttt{/012/378/456/} $\to$ \texttt{/024/138/567/} \\
$T_3$ & 9  & 1 & 91 & $(1, 4, 4), (2, 3, 4)^2$ & \texttt{/015/247/368/} \texttt{[snn]} $\to$ \texttt{/012/367/458/} $\to$ \texttt{/024/178/356/} \\
$T_4$ & 9  & 1 & 85 & $(1, 3, 5)^2, (1, 4, 4)$ & \texttt{/014/256/378/} \texttt{[nns]} $\to$ \texttt{/028/134/567/} $\to$ \texttt{/047/135/268/} \\
$T_5$ & 9  & 1 & 85 & $(2, 3, 4)^2, (3, 3, 3)$ & \texttt{/025/147/368/} \texttt{[nsn]} $\to$ \texttt{/014/258/367/} $\to$ \texttt{/028/147/356/} \\
$T_6$ & 18 & 1 & 88 & $(1, 3, 5), (1, 4, 4), (2, 3, 4)$ & \texttt{/014/267/358/} \texttt{[nsn]} $\to$ \texttt{/028/167/345/} $\to$ \texttt{/047/168/235/} \\
$T_7$ & 18 & 1 & 79 & $(1, 3, 5), (2, 2, 5), (2, 3, 4)$ & \texttt{/014/268/357/} \texttt{[nns]} $\to$ \texttt{/028/156/347/} $\to$ \texttt{/047/123/568/} \\
$T_8$ & 18 & 0 & 84 & $(1, 3, 5), (2, 3, 4)^2$ & \texttt{/014/257/368/} \texttt{[nnn]} $\to$ \texttt{/028/145/367/} $\to$ \texttt{/047/128/356/} \\
\bottomrule
\end{tabular}
\caption{\label{tab:d9-trimers}The complete catalogue of eight trimers of shape $z_3^3$ on $\mathbb{Z}_9$ ordered under the 5-tier natural sorting cascade ($|O| \uparrow$, $s \downarrow$, $\mathcal{E} \downarrow$, strides, lexicographical order)}
\end{table}

\begin{figure}[p]
\centering
\input{tkz/fig_d9_universe.tkz}
\caption{\label{fig:d9-universe}The complete $D_9$ universe: the sole Monomer ($M_1$, top) and the eight Trimers ($T_1 \dots T_8$, bottom), shown as horizontal 3-strips under the Galois generator $\mu_2 \colon x \mapsto 2x \pmod 9$}
\end{figure}
